\section*{第 \ref{ch:b3:fouriertransform} 章——
傅里叶变换}

\begin{solution}{exo:b3:fouriertransform:1}
$\widehat{\mathbf 1_{\intcc{-a}a}}(\xi) =
\int_{-a}^a\eu^{-\iu\xi x}\dd x = \frac{2\sin(a\xi)}{\xi}$
（$2a$ 为 $0$ 处的值）。
$\widehat{\eu^{-a\abs x}}(\xi) = \int_0^\infty\eu^{-(a +
\iu\xi)x} + \eu^{-(a - \iu\xi)x}\,\dd x = \frac1{a + \iu\xi} +
\frac1{a - \iu\xi} = \frac{2a}{a^2 + \xi^2}$。
帐篷函数：$\max(0, 1 - \abs x) = \mathbf 1_{\intcc{-1/2}{1/2}} *
\mathbf 1_{\intcc{-1/2}{1/2}}$，故其变换为
$\bigl(\frac{2\sin(\xi/2)}\xi\bigr)^2 =
\bigl(\frac{\sin(\xi/2)}{\xi/2}\bigr)^2$。
最后：$\frac{2a}{a^2+\xi^2} \in L^1$，故由反演公式
（\cref{thm:b3:fouriertransform:inversion}(2)）作用于
$\eu^{-a\abs x}$，变量换名后得
\[
\widehat{\Bigl(\frac1{x^2 + a^2}\Bigr)}(\xi) =
\frac{\pi}{a}\,\eu^{-a\abs\xi} .
\]
\end{solution}

\begin{solution}{exo:b3:fouriertransform:2}
由\cref{prop:b3:fouriertransform:basic}可得：
$\widehat{f(\cdot - a)} = \eu^{-\iu a\xi}\hat f(\xi)$；
$\widehat{f\cos(b\cdot)} = \frac12\bigl(\hat f(\xi - b) + \hat
f(\xi + b)\bigr)$；
$\widehat{f(a\cdot + b)}(\xi) = \frac1a\,\eu^{\iu
b\xi/a}\,\hat f(\xi/a)$（$a > 0$）；
$\widehat{\overline{f(-\cdot)}} = \overline{\hat f}$；
$\widehat{f * f} = \hat f^2$。对高斯函数
（$\widehat{\eu^{-x^2}} = \sqrt\pi\eu^{-\xi^2/4}$）逐条验证都只需一行计算——
例如 $\eu^{-(x-a)^2}$ 的变换为
$\sqrt\pi\,\eu^{-\iu a\xi}\eu^{-\xi^2/4}$，
这可由直接计算（配方）加以验证。
\end{solution}

\begin{solution}{exo:b3:fouriertransform:3}
(a) $h = \mathbf 1_{\intcc{-1}1}*\mathbf 1_{\intcc{-1}1}$ 的变换为
$\hat h = \bigl(\frac{2\sin\xi}\xi\bigr)^2 \in L^1$；
在 $x = 0$ 处应用反演，其中 $h(0) = \lambda(\intcc{-1}1\cap\intcc{-1}1)
= 2$：
\[
2 = \frac1{2\pi}\int_\R\Bigl(\frac{2\sin\xi}\xi\Bigr)^2
\dd\xi
\ \Longrightarrow\
\int_\R\Bigl(\frac{\sin\xi}\xi\Bigr)^2\dd\xi = \pi ,
\]
与 $\int_0^\infty\frac{\sin^2}{\xi^2} =
\frac\pi2$ 相符（见\cref{pb:b3:lebesgue:1}）。

(b) 取 $f = \eu^{-\abs x}$，普朗歇尔等式 $\int\abs{\hat f}^2 =
2\pi\int\abs f^2$ 即 $\int\frac{4\,\dd\xi}{(1 + \xi^2)^2} =
2\pi\int\eu^{-2\abs x}\dd x = 2\pi$：
$\int_\R\frac{\dd\xi}{(1+\xi^2)^2} = \frac\pi2$。
\end{solution}

\begin{solution}{exo:b3:fouriertransform:4}
(a) $g_t(x) = \frac1{2\sqrt{\pi t}}\eu^{-x^2/4t}$：由
\cref{ex:b3:fouriertransform:gaussian}，取 $a = \frac1{4t}$，
$\hat g_t(\xi) = \frac1{2\sqrt{\pi t}}\sqrt{4\pi
t}\,\eu^{-t\xi^2} = \eu^{-t\xi^2}$。
(b) $\widehat{g_t * g_s} = \eu^{-t\xi^2}\eu^{-s\xi^2} =
\widehat{g_{t+s}}$，且变换在 $L^1$ 上是单射
（\cref{thm:b3:fouriertransform:inversion}(3)），故 $g_t * g_s =
g_{t+s}$。
(c) $\norm{g_t}_1 = 1$（高斯积分）；
$\norm{g_t}_2^2 = \frac1{4\pi t}\int\eu^{-x^2/2t}\dd x =
\frac{\sqrt{2\pi t}}{4\pi t} = \frac1{\sqrt{8\pi t}}$。
\end{solution}

\begin{solution}{exo:b3:fouriertransform:5}
对实偶函数 $f$：$\hat f(\xi) = \int f\cos(\xi x)\dd x$（
正弦部分相消）：实值且为偶函数。奇函数：$\hat f(\xi) =
-\iu\int f\sin(\xi x)$：为奇函数且纯虚。$\hat f(0) =
\int f$：即总质量。若 $f \geq 0$：$\abs{\hat f(\xi)}
\leq \int\abs f = \int f = \hat f(0)$，故上确界在
$0$ 处取得。对于概率密度，$\hat f(-\xi)$ 即
\cref{ch:b3:clt}中的特征函数：模长 $\leq 1$，
处处成立，且在原点处 $= 1$。
\end{solution}

\begin{solution}{exo:b3:fouriertransform:6}
(i) 单射性由
\cref{thm:b3:fouriertransform:inversion}(3)给出；连续性由
$\norm{\hat f}_\infty \leq \norm f_1$ 给出（取值在
$\mathcal C_0$ 中，由黎曼–勒贝格引理）；$\mathcal C_0$ 在
上确界范数下是闭的（无穷远处趋于零的函数的一致极限，
仍在无穷远处趋于零）：故为巴拿赫空间。

(ii) 巴拿赫空间之间的连续双射，其
逆映射连续（\cref{thm:b3:banach:openmapping}）：存在
常数 $C$，使得 $\norm f_1 \leq C\norm{\hat f}_\infty$。

(iii) 令 $f_n(x) = \frac{\sin x}x\cdot\frac{\sin(x/n)}{x/n}$：
两个 $L^2$ 函数之积，且
在无穷远处为 $O(x^{-2})$，故 $f_n \in L^1\cap L^2$。由于
$\bigl(\frac{\sin(ax)}{ax}\bigr)$ 的 $L^2$ 变换为
$\frac\pi a\mathbf 1_{\intcc{-a}a}$，乘积公式
$\widehat{gh} = \frac1{2\pi}\hat g * \hat h$（对 $g, h
\in L^2$ 且 $gh \in L^1$ 成立；在施瓦茨函数上用富比尼验证，
再借两端关于 $L^2$ 的连续性推广，依普朗歇尔
定理）给出
\[
\hat f_n = \frac1{2\pi}\,\bigl(\pi\mathbf
1_{\intcc{-1}1}\bigr) * \bigl(\pi n\,\mathbf
1_{\intcc{-1/n}{1/n}}\bigr):
\]
这是一个高为 $\frac{\pi n}2\cdot\frac2n = \pi$ 的梯形：
$\norm{\hat f_n}_\infty = \pi$ 对每个 $n$ 成立。但在 $[1,
n]$ 上，$\frac{\sin(x/n)}{x/n} \geq \sin 1 > 0$，故
\[
\norm{f_n}_1 \geq \sin 1\int_1^n\frac{\abs{\sin x}}x\dd x
\geq c\ln n
\]
（拱计数法，见\cref{thm:b3:banach:fourierdiverge}）。
界 $\norm{f_n}_1 \leq C\pi$ 对大的 $n$ 不成立，故该映射不是
满射。（其像是一个稠密的——由
斯通–魏尔斯特拉斯型论证得到——但为
$\mathcal C_0$ 的真子空间。）
\end{solution}

\begin{solution}{exo:b3:fouriertransform:7}
若 $\hat f(\xi) = O(\abs\xi^{-k-1-\delta})$，则 $\xi^j\hat f
\in L^1$ 对 $0 \leq j \leq k$ 成立（无穷远处的\omterm{def:b3:lebesgue:l1}{可积}性由
衰减性推出，局部则由 $\hat f$ 的连续性得到）。反演
（\cref{thm:b3:fouriertransform:inversion}(2)）表示 $f$，
几乎处处由 $x \mapsto \frac1{2\pi}\int\hat f(\xi)\eu^{\iu
x\xi}\dd\xi$ 给出，积分号下的求导
（控制函数为 $\abs{\xi^j\hat f}$）使该代表元
属于 $\mathcal C^k$。反之，对 $f \in \mathcal C_c^k$：
反复应用\cref{prop:b3:fouriertransform:basic}(4)，
$(\iu\xi)^k\hat f = \widehat{f^{(k)}}$，故 $\abs{\hat f} \leq
\norm{f^{(k)}}_1\abs\xi^{-k}$。帐篷函数连续，
且具紧支集（$k = 0$：变换有界），其
变换 $\sim \xi^{-2} = O(\abs\xi^{-0-1-1})$ 由第一个方向，
给出一个 $\mathcal C^0$ 的代表元——两者都是
最优的：帐篷函数不属于 $\mathcal C^1$，其变换
的衰减不快于 $\xi^{-2}$。
\end{solution}

\begin{solution}{exo:b3:fouriertransform:8}
分部积分（$f \in \mathcal S$；边界项
均为零）：
\[
1 = \int f^2 = \bigl[xf^2\bigr]_{-\infty}^{\infty} - \int
x\,(f^2)' = -2\int xff' \leq 2\,\norm{xf}_2\,\norm{f'}_2 .
\]
由普朗歇尔恒等式及 $\widehat{f'} = \iu\xi\hat f$：
$\norm{f'}_2^2 = \frac1{2\pi}\int\xi^2\abs{\hat f}^2$。将
上式平方得：
\[
\frac14 \leq \norm{xf}_2^2\cdot\frac1{2\pi}
\int\xi^2\abs{\hat f}^2\dd\xi .
\]
等号成立要求柯西–施瓦茨不等式取等，即 $f' =
\lambda xf$，其中 $\lambda < 0$（由\omterm{def:b3:lebesgue:l1}{可积}性），亦即 $f(x) =
c\,\eu^{\lambda x^2/2}$：高斯函数。在
$x$ 中集中（$\norm{xf}_2$ 小）的信号，其频谱必须分散，且
反之亦然，此即不确定性原理。
\end{solution}

\begin{solution}{exo:b3:fouriertransform:9}
高斯函数 $f(x) = \eu^{-\pi tx^2}$ 属于\omterm{def:b3:fouriertransform:schwartz}{施瓦茨函数类}，故
\cref{thm:b3:fouriertransform:poisson} 适用，并且
$\hat f(\xi) = t^{-1/2}\eu^{-\xi^2/4\pi t}$。取 $\xi = 2\pi
k$ 时右端化为 $t^{-1/2}\eu^{-\pi k^2/t}$，即
$\theta$ 恒等式。取 $t = 1$：
\[
\sum_{n\in\Z}\eu^{-\pi n^2} = 1 + 2\eu^{-\pi} + 2\eu^{-4\pi} +
\cdots \approx 1 + 0.0864278 + 0.0000070 = 1.0864348,
\]
取三项即精确到小数点后 $6$ 位（$\eu^{-9\pi}
\approx 5\cdot10^{-13}$）。当 $t = 10^{-2}$ 时，定义所给的
级数需要 $\eu^{-\pi n^2/100} < 10^{-7}$，即 $n \gtrsim
23$——约 $47$ 项——而变换后的级数为
$10\sum_k\eu^{-100\pi k^2}$，其中 $k = 1$ 项
已为 $\sim 10^{-136}$，只取一项即可。
\end{solution}

\begin{solution}{exo:b3:fouriertransform:10}
$\mathcal Ff \in L^2(\intcc{-\pi}\pi) \subseteq
L^1(\intcc{-\pi}\pi)$（有限\omterm{def:b3:measure:measure}{测度}），故反演给出
连续代表元 $f(x) =
\frac1{2\pi}\int_{-\pi}^\pi\mathcal Ff(\xi)\eu^{\iu
x\xi}\dd\xi$，其中
\[
f(n) = \frac1{2\pi}\int_{-\pi}^{\pi}\mathcal
Ff(\xi)\,\eu^{\iu n\xi}\dd\xi = \langle e_{-n}, \mathcal
Ff\rangle
\]
按\cref{thm:b3:hilbert:fourier}的记号，
在该\omterm{def:b3:hilbert:onb}{希尔伯特基}下展开：$\mathcal Ff = \sum_nf(n)\,\eu^{-\iu
n\xi}$ 在 $L^2(\intcc{-\pi}\pi)$ 中。
对 $L^2$ 连续的 $\mathcal F^{-1}$ 逐项应用：
\[
\mathcal F^{-1}\bigl(\mathbf
1_{\intcc{-\pi}\pi}\eu^{-\iu n\xi}\bigr)(x)
= \frac1{2\pi}\int_{-\pi}^{\pi}\eu^{\iu\xi(x - n)}\dd\xi
= \frac{\sin\bigl(\pi(x-n)\bigr)}{\pi(x - n)} ,
\]
于是得到 $f = \sum_nf(n)\operatorname{sinc}(\cdot - n)$，在
$L^2$ 意义下：带限信号由其整数
采样值唯一确定——即香农采样定理。
\end{solution}

\begin{solution}{exo:b3:fouriertransform:11}
(a) $\mathcal S$ 上的反演：$f(x) =
\frac1{2\pi}\int\hat f(\xi)\eu^{\iu x\xi}\dd\xi =
\frac1{2\pi}(\mathcal F\hat f)(-x)$，即 $(\mathcal
F^2f)(x) = 2\pi f(-x)$。两次应用得 $\mathcal F^4f =
2\pi\,\mathcal F^2f(-\cdot) = (2\pi)^2f$.

(b) 若 $\mathcal Ff = \lambda f$ 且 $f \neq 0$：
$(2\pi)^2f = \mathcal F^4f = \lambda^4f$，故 $\lambda^4 =
(2\pi)^2$：$\lambda \in \{\pm\sqrt{2\pi},
\pm\iu\sqrt{2\pi}\}$。高斯函数 $g(x) = \eu^{-x^2/2}$ 满足
$\hat g = \sqrt{2\pi}\,g$
（\cref{ex:b3:fouriertransform:gaussian} 中取 $a = \frac12$）：
以 $+\sqrt{2\pi}$ 为特征值的特征函数。

(c) $\mathcal F^2f = 2\pi f(-\cdot)$ 等于 $\pm2\pi f$，
取决于奇偶性。令 $h(x) = x\eu^{-x^2/2}$，
对 $\hat g(\xi) = \sqrt{2\pi}\eu^{-\xi^2/2}$ 求导，
利用 $\widehat{xf} = \iu\frac{\dd}{\dd\xi}\hat f$：
\[
\hat h(\xi) = \iu\,\frac{\dd}{\dd\xi}\bigl(\sqrt{2\pi}
\eu^{-\xi^2/2}\bigr) = -\iu\sqrt{2\pi}\,\xi\eu^{-\xi^2/2}
= -\iu\sqrt{2\pi}\,h(\xi) :
\]
以 $-\iu\sqrt{2\pi}$ 为特征值的特征函数。（埃尔米特
函数延续这一模式，周而复始地经过四个
特征值——构成离散傅里叶“时钟”。）

\end{solution}

\begin{solution}{exo:b3:fouriertransform:12}
(a) $\tilde f \in L^2$，且 $\norm{\tilde f}_2 = \norm f_2$；
由 \cref{exo:b3:lp:6}（共轭指数 $p = q = 2$）
可知 $A_f = f * \tilde f$ 有界且一致连续，并有
\[
A_f(x) = \int f(y)\,\overline{f(y - x)}\,\dd y,
\qquad A_f(0) = \norm f_2^2,
\qquad \abs{A_f(x)} \leq \norm f_2\,\norm{f(\cdot -
x)}_2 = A_f(0)
\]
由柯西--施瓦茨不等式。

(b) 对 $f \in L^1\cap L^2$，$\tilde f \in L^1$ 同样成立，并由
卷积定理给出 $\widehat{A_f} = \hat f\,
\widehat{\tilde f}$；计算得 $\widehat{\tilde f}(\xi) =
\int\overline{f(-x)}\eu^{-\iu\xi x}\dd x =
\overline{\int f(u)\eu^{-\iu\xi u}\dd u} =
\overline{\hat f(\xi)}$：$\widehat{A_f} = \abs{\hat f}^2
\geq 0$。

(c) 当 $\abs{\hat f}^2 \in L^1$ 时，反演公式可用于
连续函数 $A_f$（其变换\omterm{def:b3:lebesgue:l1}{可积}；
\cref{thm:b3:fouriertransform:inversion}）：
\[
A_f(x) = \frac1{2\pi}\int\abs{\hat f(\xi)}^2
\eu^{\iu x\xi}\,\dd\xi .
\]
对 $f = \mathbf 1_{\intcc{-1/2}{1/2}}$，有 $\hat f(\xi) =
\frac{2\sin(\xi/2)}\xi = \frac{\sin(\xi/2)}{\xi/2}$，并在
$x = 0$ 处：
\[
1 = \norm f_2^2 = \frac1{2\pi}\int_\R
\Bigl(\frac{\sin(\xi/2)}{\xi/2}\Bigr)^2\dd\xi
= \frac1{2\pi}\cdot2\int_\R\Bigl(\frac{\sin u}u\Bigr)^2\dd u
\]
（$\xi = 2u$），即 $\int_\R\bigl(\frac{\sin u}u\bigr)^2\dd
u = \pi$——普朗歇雷尔最喜欢的积分，借助
自相关求得。
\end{solution}

\begin{solution}{pb:b3:fouriertransform:1}
\textbf{1.} 对方程关于 $x$ 形式地作变换：
$\partial_t\hat u(t,\xi) = \widehat{\partial^2_{xx}u} =
(\iu\xi)^2\hat u = -\xi^2\hat u$，这是关于 $t$ 的常微分方程，对每个
频率，解为 $\hat u(t,\xi) = \eu^{-t\xi^2}\hat f(\xi)$。由于
$\eu^{-t\xi^2} = \hat g_t$
（\cref{exo:b3:fouriertransform:4}），故此乘积为
$g_t * f$ 的变换。

\textbf{2.} $\abs{u(t,x)} \leq \norm f_\infty\int g_t = \norm
f_\infty$：良定义。在 $[t_0, T]\times[-A, A]$ 上，每个
混合偏导数 $\partial^m_t\partial^n_xg_t(x - y)$ 是
关于 $(x - y)$ 与 $t^{-1}$ 的多项式，乘以
$\eu^{-(x-y)^2/4t}$，当 $\abs y \geq 2A$ 时，以
$C\,(1 + y^2)^N\eu^{-(\abs y - A)^2/4T}$ 为界，是一个\omterm{def:b3:lebesgue:l1}{可积}的
与窗口内 $(t, x)$ 无关的控制函数（且
当 $\abs y \leq 2A$ 时有界）：对积分号下的
重复求导（\cref{thm:b3:lebesgue:paramdiff}）可用：$u \in
\mathcal C^\infty(\intoo0\infty\times\R)$。

\textbf{3.} 记 $g_t(x) = \frac1{2\sqrt{\pi
t}}\eu^{-x^2/4t}$：
\[
\partial_tg_t = g_t\Bigl(\frac{x^2}{4t^2} - \frac1{2t}\Bigr)
= \partial^2_{xx}g_t
\]
（对 $x$ 求两次导：$\partial_xg_t = -\frac
x{2t}g_t$，$\partial^2_{xx}g_t = \bigl(\frac{x^2}{4t^2} -
\frac1{2t}\bigr)g_t$）。由第 2 问，导数可
移入积分号下：$\partial_tu = \partial^2_{xx}u$。

\textbf{4.} $u(t,x) - f(x) = \int g_t(y)\bigl(f(x - y) -
f(x)\bigr)\dd y$。给定紧集 $K$ 与 $\varepsilon$：
由 $f$ 在 $K$ 的某\omterm{def:b3:topology:topology}{邻域}上一致连续，可得
存在 $\delta$，使得 $\abs{f(x-y) - f(x)} < \varepsilon$ 对 $x \in
K$，$\abs y \leq \delta$；尾部的贡献 $\leq 2\norm
f_\infty\int_{\abs y > \delta}g_t(y)\dd y = 2\norm
f_\infty\,\P$-质量在 $\delta$ 之外，其值为
$\frac2{\sqrt\pi}\int_{\delta/2\sqrt t}^\infty\eu^{-z^2}\dd z
\to 0$，当 $t \to 0$ 时。对 $f \in L^p$：$\norm{u(t) - f}_p \leq
\int g_t(y)\norm{\tau_yf - f}_p\dd y$（闵可夫斯基与琴生，参见
\cref{thm:b3:lp:regularization}），同样分段，利用
\cref{thm:b3:lp:density}(3)。

\textbf{5.} 瞬时光滑化即第 2 问：$u(t)$ 是
$\mathcal C^\infty$（$t > 0$），且未用到 $f$ 的
光滑性。对 $f = \mathbf 1_{\intoo0\infty}$：
\[
u(t, x) = \int_0^\infty g_t(x - y)\dd y
= \frac1{\sqrt\pi}\int_{-x/2\sqrt t}^{\infty}\eu^{-z^2}\dd z
= \frac12\Bigl(1 +
\operatorname{erf}\Bigl(\frac{x}{2\sqrt t}\Bigr)\Bigr),
\qquad \operatorname{erf}(s) =
\frac2{\sqrt\pi}\int_0^s\eu^{-z^2}\dd z :
\]
光滑化阶跃，其过渡区以 $\sqrt t$ 的量级变宽
（$t_1 < t_2 < t_3$ 时的图像：越来越平缓的斜坡经过
$(0, \frac12)$）。

\textbf{6.} 被积函数 $g_t(x-y)f(y)$ 满足 $\geq 0$，而
核严格为正：$u(t,x) = 0$ 则必有 $f = 0$
几乎处处。关于 $f$ 的单调性即积分的单调性。
即便 $f = 0$ 在某区间上，$u(t, \cdot) > 0$ 在该处
对每个 $t > 0$ 均成立：热量以无穷大速度传播，
（任何地方的正性都瞬时被各处感受到）。

\textbf{7.} 由托内利定理（$g_t(x-y)\abs{f(y)}$ 在
$\R^2$ 上\omterm{def:b3:lebesgue:l1}{可积}）：$\int u(t,x)\dd x = \int f(y)\bigl(\int g_t(x -
y)\dd x\bigr)\dd y = \int f$。

\textbf{8.} 普朗歇尔等式：$2\pi\norm{u(t)}_2^2 =
\int\eu^{-2t\xi^2}\abs{\hat f(\xi)}^2\dd\xi$，关于
$t$ 逐点单调不增，且严格递减，除非 $\hat f = 0$ 几乎处处（即 $= f =
0$）；由控制收敛定理，极限为 $0$（$t\to\infty$ 时）。此外
$\norm{u(t)}_\infty \leq \norm{g_t}_\infty\norm f_1 =
\frac{\norm f_1}{2\sqrt{\pi t}} \to 0$。

\textbf{9.} 对 $\varphi \in \mathcal C_c^\infty$，$t \mapsto
\langle\varphi, \hat u(t)\rangle$ 为 $\mathcal C^1$ 函数，其
导数为 $\langle\varphi, \partial_t\hat u\rangle =
\langle\varphi, \widehat{\partial_{xx}u}\rangle =
\langle\xi^2\varphi\dots\rangle$——准确地说，
$\widehat{\partial^2_{xx}u} = -\xi^2\hat u$ 将原方程
转化为变换域上的方程。于是对几乎处处的 $\xi$，绝对连续
函数 $t \mapsto \eu^{t\xi^2}\hat u(t,\xi)$ 的导数为
$\eu^{t\xi^2}(\xi^2\hat u + \partial_t\hat u) = 0$，在
积分意义下，它为常数，并令 $t \to 0$
（$\hat u(t) \to \hat f$ 在 $L^2$ 中，且沿某子列几乎处处）：
$\hat u(t, \xi) = \eu^{-t\xi^2}\hat f(\xi)$ 几乎处处成立。该类中的两个
解具有相同的变换，故二者
相等。

\textbf{10.} 若 $f = g_s * h$ 且 $h \in L^2$，则必有 $\hat f =
\eu^{-s\xi^2}\hat h$，即 $\hat h = \eu^{s\xi^2}\hat f \in
L^2$。取 $\hat f(\xi) = \eu^{-\abs\xi}$，则 $f(x) =
\frac1\pi\cdot\frac1{1 + x^2}$
（见\cref{exo:b3:fouriertransform:1}，由反演得到），这是一个
非常光滑的 $L^2$ 函数；但 $\eu^{2s\xi^2 - 2\abs\xi} \to
\infty$：$\eu^{s\xi^2}\hat f \notin L^2$，对每个 $s > 0$ 均成立。
柯西型函数\emph{绝不}是先前
扩散的结果。

\textbf{11.} 热半群将变换乘以
$\eu^{-t\xi^2}$，该因子处处不为零，故为单射——
形式上没有信息丢失。但其值域由
变换按 $\eu^{-t\xi^2}$ 速率衰减的函数组成，是
$L^2$ 中一个很小、稠密而真的子空间（第 10 题表明
即使是很好的函数也在其外）。求逆将
把频率 $\xi$ 放大 $\eu^{t\xi^2}$ 倍：无界，
因而对任何扰动都不稳定。扩散
之所以不可逆，并非因为映射遗忘信息，而是因为其
逆不能连续——由
泛函分析写成的时间之箭。

\textbf{12.} $\hat f \in L^2(\intcc{-\Omega}\Omega)
\subseteq L^1$（由柯西--施瓦茨不等式，区间有界），故
$F(x) = \frac1{2\pi}\int_{-\Omega}^\Omega\hat
f(\xi)\eu^{\iu x\xi}\dd\xi$ 处处有定义，且
积分号下求导（以该带上的
$\Omega^k\abs{\hat f} \in L^1$ 为控制）使其
成为 $\mathcal C^\infty$ 函数，且 $\abs{F^{(k)}} \leq
\frac{\Omega^k}{2\pi}\norm{\hat f}_{L^1}$ 处处成立。又 $F
= f$ 几乎处处：两边的变换相同，而
该变换在 $L^2$ 上为单射
（见\cref{thm:b3:fouriertransform:plancherel} 及其
$L^2$ 上的推广）。

\textbf{13.} 指数函数族 $\xi \mapsto \eu^{-\iu
n\pi\xi/\Omega}$，$n \in \Z$，构成
$L^2(\intcc{-\Omega}\Omega)$
的\omterm{def:b3:hilbert:onb}{希尔伯特基}（由\cref{thm:b3:hilbert:fourier}经尺度变换得到）。其
$\hat f$ 沿第 $n$ 个基元的系数为
\[
\frac1{2\Omega}\int_{-\Omega}^\Omega\hat f(\xi)\,
\eu^{\iu n\pi\xi/\Omega}\dd\xi
= \frac{2\pi}{2\Omega}\cdot
\frac1{2\pi}\int_{-\Omega}^{\Omega}\hat f(\xi)\,
\eu^{\iu(n\pi/\Omega)\xi}\dd\xi
= \frac\pi\Omega\,f\Bigl(\frac{n\pi}\Omega\Bigr),
\]
由第 12 题的公式，取 $x = \frac{n\pi}\Omega$，
所述展开在该带上的 $L^2$ 意义下成立。

\textbf{14.} 将展开式代入
第 12 题的反演公式；求和与积分的交换即为
$L^2$ 配对关于 $\frac1{2\pi}
\eu^{\iu x\xi}\mathbf 1_{\abs\xi\leq\Omega}$ 的连续性（该函数的 $L^2$
范数 $\frac{\sqrt{2\Omega}}{2\pi}$，与 $x$ 无关——
由此得一致性）：
\[
f(x) = \sum_nf\Bigl(\frac{n\pi}\Omega\Bigr)\cdot
\frac1{2\Omega}\int_{-\Omega}^\Omega
\eu^{\iu(x - n\pi/\Omega)\xi}\dd\xi
= \sum_nf\Bigl(\frac{n\pi}\Omega\Bigr)
\operatorname{sinc}(\Omega x - n\pi),
\]
因为 $\frac1{2\Omega}\int_{-\Omega}^\Omega\eu^{\iu
u\xi}\dd\xi = \frac{\sin(\Omega u)}{\Omega u}$.

\textbf{15.} 将第 14 题的计算倒过来看，
$s_n = \operatorname{sinc}(\Omega\cdot -
n\pi)$ 的变换为 $\hat s_n = \frac\pi\Omega\,\eu^{-\iu
n\pi\xi/\Omega}\,\mathbf 1_{\intcc{-\Omega}\Omega}$。
普朗歇尔等式：
\[
\langle s_n, s_m\rangle = \frac1{2\pi}
\Bigl(\frac\pi\Omega\Bigr)^2\int_{-\Omega}^\Omega
\eu^{\iu(n-m)\pi\xi/\Omega}\dd\xi =
\frac\pi\Omega\,\delta_{nm} :
\]
这是一族具有常数范数的正交族，
其范数为 $\sqrt{\pi/\Omega}$。对第 14 题的
展开式两边取范数：$\norm f_2^2 =
\frac\pi\Omega\sum_n\abs{f(n\pi/\Omega)}^2$.

\textbf{16.} $g(x) = \sin(\Omega x)\operatorname{sinc}
(\Omega x) = \frac{\sin^2(\Omega x)}{\Omega x}$ 在每个格点
$\frac{n\pi}\Omega$ 处为零（包括 $0$，由极限得到），
且不恒为零。其频带：写 $g =
\frac1{2\iu}\bigl(\eu^{\iu\Omega x} - \eu^{-\iu\Omega
x}\bigr)\operatorname{sinc}(\Omega x)$；
以 $\eu^{\pm\iu\Omega x}$ 作调制，使变换平移
$\mp\Omega$，故 $\hat g$ 支撑于
$\intcc{-2\Omega}{2\Omega}$（实际上在两个
平移带之并中）：$g \in PW_{2\Omega}$，对
$\Omega$ 速率采样不可见——混叠的化身。

\textbf{17.} 由第 15 题，样本
平等地承载能量：$\norm f^2 = \frac\pi\Omega\sum
\abs{f(n\pi/\Omega)}^2$。若信号在
时间窗口 $\intcc{-T/2}{T/2}$ 之外的能量为 $\leq \varepsilon^2\norm
f^2$，则窗外的样本满足（至多相差
由第 12 题的一致界所控制的边界项）：
$\frac\pi\Omega\sum_{\abs{n\pi/\Omega} > T/2}
\abs{f(n\pi/\Omega)}^2 \approx \norm{f\,\mathbf 1_{\abs x >
T/2}}^2 \leq \varepsilon^2\norm f^2$：将采样级数
截断至 $\approx \frac{\Omega T}\pi$
个窗内指标，即可将 $f$ 重构至相对误差
$\approx\varepsilon$。因此时间--带宽积
$\frac{\Omega T}{\pi}$ 计数了信号的有效实
自由度——这正是每一种音频
格式背后的规则。

\textbf{18.} (a) $\operatorname{sinc}(\Omega x)$ 的
样本为 $f(n\pi/\Omega) = \operatorname{sinc}(n\pi) =
\delta_{n0}$：级数只剩下其 $n = 0$ 项，
即 $\operatorname{sinc}(\Omega x)$——定理再现
其自身的核。(b) 若 $\hat f$ 支撑于
$\intcc{-\Omega'}{\Omega'} \subseteq
\intcc{-\Omega}\Omega$，第 13--14 题的每一步
原样成立，取较大的带宽 $\Omega$（
$\hat f$ 在较大区间上的展开仍然合法）：
以快于自身奈奎斯特速率的频率采样，
对重构毫无影响——过采样无害，且在
实践中有益（此时可采用
衰减更快的重构核）。

\textbf{19.} 在积分内展开 $\eu^{-\iu\xi x} =
\sum_k\frac{(-\iu\xi x)^k}{k!}$；在
$\intcc{-A}A$ 上级数正规收敛
（$\sum_k\frac{\abs{\xi}^kA^k}{k!}\abs f \in L^1$），故
逐项积分是合法的：
\[
\hat f(\xi) = \sum_{k\geq0}\frac{(-\iu\xi)^k}{k!}m_k,
\qquad \abs{m_k} \leq A^k\norm f_1 .
\]
此界保证该级数对每个复数 $\xi$ 收敛；
在任一点 $\xi_0$ 附近，重新组合各项（由绝对收敛），
给出关于 $\xi - \xi_0$ 的幂级数：$\hat f$ 是
实解析的，且各点处收敛半径皆为无穷。

\textbf{20.} 设 $g$ 为 $\R$ 上的实解析函数（泰勒级数
在各点附近收敛于 $g$），并记 $Z = \{\xi :
g^{(k)}(\xi) = 0\ \forall k\}$。$Z$ 为闭集（
闭集之交）；它也是开的，因为在 $\xi_0 \in Z$ 处的局部
泰勒展开（$g$ 的）为零级数，故 $g$
在 $\xi_0$ 附近恒为零，连同其全部导数。若
$g$ 在某区间上为零，则 $Z \neq \varnothing$；由
$\R$ 的连通性，$Z = \R$：$g \equiv 0$。若 $f
\neq 0$ 与 $\hat f$ 皆具紧支集：由
第 19 题使 $\hat f$ 为实解析，在某紧集外为零，
从而在区间上为零：$\hat f \equiv 0$，故由
单射性得 $f = 0$ 几乎处处——矛盾。同理，非零的 $f \in
PW_\Omega$ 不可能是紧支的（借助反演交换
$f$ 与 $\hat f$ 的角色）：带限信号永不消失，
时限信号则占据无界谱。

\textbf{21.} 对 $f = \eu^{-ax^2}$：$\norm f_2^2 =
\sqrt{\frac\pi{2a}}$，且 $\int x^2\abs f^2 =
\frac1{4a}\sqrt{\frac{\pi}{2a}}$（高斯二阶矩）；
$\hat f = \sqrt{\frac\pi a}\,\eu^{-\xi^2/4a}$
（见\cref{ex:b3:fouriertransform:gaussian}），且
\[
\frac1{2\pi}\int\xi^2\abs{\hat f}^2\dd\xi =
\frac1{2\pi}\cdot\frac\pi a\int\xi^2
\eu^{-\xi^2/2a}\dd\xi = \frac1{2a}\cdot a\sqrt{2\pi a}
= \frac{\sqrt{2\pi a}}2 .
\]
归一化之积：$\frac1{4a}\sqrt{\frac\pi{2a}}\cdot
\frac{\sqrt{2\pi a}}2\big/\frac{\pi}{2a} = \frac14$,
与 $a$ 无关。缩放解释了这一恒定性：将
$f$ 换为 $f(\lambda\cdot)$ 会使 $\int x^2\abs f^2/\norm
f^2$ 乘以 $\lambda^{-2}$，而 $\frac1{2\pi}\int\xi^2\abs{\hat
f}^2/\norm f^2$ 乘以 $\lambda^{2}$：该乘积为伸缩
不变量，且高斯函数构成一条伸缩\omterm{def:b3:groups:action}{轨道}。

\textbf{22.} 以 $\abs f^2$ 为位置密度，并以
$\frac1{2\pi}\abs{\hat f}^2$ 为某
量子态的动量密度（物理单位中插入 $\hbar$），
由\cref{exo:b3:fouriertransform:8}可知 $\sigma_x\sigma_p
\geq \frac\hbar2$：没有态能在两个可观测量上同时尖锐。
贯穿全章，同一条规律有五种面貌：热
瞬时抹平，因为 $\eu^{-t\xi^2}$ 消去了高频
（第二部分）；流不能倒转，因为恢复
它们是无界的（第四部分）；带限信号
足够刚性，可以落在可数格点上（第五部分）；没有函数
能突破海森堡下界；且没有函数在
变换两侧均具紧支集（第 19--20 题）。
$\hat f$ 在无穷远处的表现，决定了 $f$
在任何地方可能具有的性质。

\textbf{23.} $g_t$ 与 $g_s$ 皆属于 $L^1$，且
$\widehat{g_t}(\xi) = \eu^{-t\xi^2}$（第 1 题的
计算），故由卷积定理得
$\widehat{g_t * g_s} = \eu^{-t\xi^2}\eu^{-s\xi^2} =
\eu^{-(t+s)\xi^2} = \widehat{g_{t+s}}$；两个 $L^1$ 函数
具有相同变换者几乎处处相等（由单射性，经由
反演定理——此处两侧皆连续，故
二者处处相等）：$g_t * g_s = g_{t+s}$。因此
$u(t + s) = g_{t+s} * f = g_s * (g_t * f) = g_s * u(t)$
（卷积的结合律，托内利定理）。对数凸性：设
$N(t) = \norm{u(t)}_2^2 = \frac1{2\pi}\int
\eu^{-2t\xi^2}\abs{\hat f}^2\dd\xi$（普朗歇尔恒等式，第
8）。对 $t = \frac{t_1 + t_2}2$，写
\[
\eu^{-2t\xi^2}\abs{\hat f}^2
= \Bigl(\eu^{-2t_1\xi^2}\abs{\hat f}^2\Bigr)^{1/2}
\Bigl(\eu^{-2t_2\xi^2}\abs{\hat f}^2\Bigr)^{1/2},
\]
由柯西–施瓦茨不等式得 $N\bigl(\frac{t_1+t_2}2\bigr) \leq
\sqrt{N(t_1)\,N(t_2)}$：$\ln N$ 为中点凸的，且
由 $t$ 的控制收敛知其连续，故为凸的；
于是 $\ln\norm{u(t)}_2 = \frac12\ln N(t)$ 亦然。衰减伴随凸的
对数：热流不能在骤降中损失能量，并
随后停滞。

\textbf{24.} 核的矩：$\int g_t = 1$（取第 7 题中的 $f = g_s$，
或直接由高斯积分），$\int
x\,g_t(x)\dd x = 0$（被积函数为奇函数），并作代换 $x =
2\sqrt t\,v$，
\[
\int_\R x^2g_t(x)\,\dd x
= \frac{4t}{\sqrt\pi}\int_\R v^2\eu^{-v^2}\dd v = 2t .
\]
在卷积中作代换 $x = z + y$，并注意
$\iint(\abs z + \abs y)^2g_t(z)f(y)\,\dd z\,\dd y < \infty$
（$\int\abs z^kg_t$ 与 $\int\abs y^kf$ 均对 $k
\leq 2$，利用 $\abs y \leq \frac{1 + y^2}2$），富比尼
定理与托内利定理可用于下列矩积分：
\[
\int x\,u(t,x)\dd x = \iint (z + y)\,g_t(z)f(y)\,\dd z\,\dd
y = 0\cdot\!\int\! f + 1\cdot\!\int\! yf(y)\dd y,
\]
即第一个结论；另有
\[
\iint (z+y)^2g_t(z)f(y)\,\dd z\,\dd y
= 2t\int f + 2\cdot0\cdot\!\int\! yf + \int y^2f(y)\dd y ,
\]
即第二个结论。均值可加，方差可加，且核
的均值为 $0$、方差为 $2t$：经过时间 $t$ 后
热量已扩散开，宽度量级为 $\sqrt{2t}$——
距离随时间的平方根增长，这正是
扩散的特征（布朗路径也具有此性质，
见\cref{ch:b3:probability}）。

\textbf{25.} 变换一侧：$\hat f(\xi) =
\sqrt\pi\,\eu^{-\xi^2/4}$，于是 $\hat u(t,\xi) =
\sqrt\pi\,\eu^{-(t + \frac14)\xi^2}$，它是
$(1 + 4t)^{-1/2}\exp\bigl(-x^2/(1+4t)\bigr)$ 的变换（由
高斯变换表 $\eu^{-ax^2} \mapsto
\sqrt{\pi/a}\,\eu^{-\xi^2/4a}$，其中 $a = \frac1{1+4t}$）：
即闭形式。直接验算，记 $\sigma = 1 + 4t$：
\[
\partial_tu = \sigma^{-1/2}\eu^{-x^2/\sigma}
\Bigl(-\frac2\sigma + \frac{4x^2}{\sigma^2}\Bigr)
= \partial^2_{xx}u ,
\]
两边均由 $\partial_xu =
-\frac{2x}\sigma\,u$ 算得。守恒：$\int u(t) =
\sigma^{-1/2}\sqrt{\pi\sigma} = \sqrt\pi$ 对一切 $t$ 成立。
耗散：
\[
\norm{u(t)}_2^2 = \frac1\sigma\int\eu^{-2x^2/\sigma}\dd x
= \frac1\sigma\sqrt{\frac{\pi\sigma}2}
= \sqrt{\frac\pi2}\,(1+4t)^{-1/2},
\]
于是 $\norm{u(t)}_2 = (\pi/2)^{1/4}(1+4t)^{-1/4}$，
非增，对数为凸（第 23 题）；
$L^2$ 范数以 $t^{-1/4}$ 的速率衰减，恰为
$t^{-1/2}$ 指数的一半（该指数即 $\norm{u(t)}_\infty$ 的衰减指数）——这与
$\norm u_2^2 \leq \norm u_\infty\norm u_1$ 以及
$\norm u_1$ 的守恒性相一致。方差：$\int x^2u(t) =
\frac1\sigma\cdot\frac{\sigma^{3/2}\sqrt\pi}2 =
\frac{\sqrt\pi}2\,(1 + 4t) = \int x^2f + 2t\sqrt\pi$，与
第 24 题的预测相符（$\int x^2f = \frac{\sqrt\pi}2$，$\int
f = \sqrt\pi$）。取 $t = 6$，$\sigma = 25$，峰值高度
为 $u(6, 0) = \frac15$（对比 $u(0,0) = 1$），宽度尺度
为初值宽度的 $\sqrt\sigma = 5$ 倍，而 $\int u =
\sqrt\pi \approx 1.7725$ 始终成立：光斑
高度降为原来的五分之一，宽度增为原来的五倍，热量分毫未失。
\end{solution}
